\documentclass{article}

% if you need to pass options to natbib, use, e.g.:
%     \PassOptionsToPackage{numbers, compress}{natbib}
% before loading neurips_2025

% ready for submission
\usepackage{neurips_2025}

% to compile a preprint version, e.g., for submission to arXiv, add add the
% [preprint] option:
%     \usepackage[preprint]{neurips_2025}

\usepackage[utf8]{inputenc} % allow utf-8 input
\usepackage[T1]{fontenc}    % use 8-bit T1 fonts
\usepackage{hyperref}       % hyperlinks
\usepackage{url}            % simple URL typesetting
\usepackage{booktabs}       % professional-quality tables
\usepackage{amsfonts}       % blackboard math symbols
\usepackage{nicefrac}       % compact symbols for 1/2, etc.
\usepackage{microtype}      % microtypography
\usepackage{xcolor}         % colors
\usepackage{graphicx}
\usepackage{subcaption}
\usepackage{amsmath,amssymb}

\begin{document}

\section*{Experiment}

This section presents the experimental evaluation of the proposed Error-Compensated Proximal Random Reshuffling (EC-PRR) method and its comparison with the baseline Proximal Gradient Random Reshuffling (PG-RR). The experiments are designed to assess convergence speed, robustness to observation and operator inexactness, per-iteration variability induced by random reshuffling, and stability of accumulated gradient trajectories in realistic finite-sample settings.

\subsection*{Experimental Setup}

\textbf{Dataset and preprocessing} The experiments use grayscale images resized to $28\times28$ and vectorized to form ground-truth signals $x^\star\in\mathbb{R}^{784}$. When image files are unavailable, a reproducible synthetic Gaussian vector is used. Pixel intensities are normalized and the initial iterate is generated by adding a small Gaussian perturbation to $x^\star$ to emulate a realistic starting point.

\textbf{Linear measurement model} For each trial we generate $m=10$ independent Gaussian measurement blocks $A_i\in\mathbb{R}^{d\times D}$ with $d=100$ and $D=784$. Noiseless observations are $y_i=A_i x^\star$. The finite-sample regime implemented corresponds to random reshuffling: each epoch performs one pass over a permutation of the $m$ measurement blocks.

\textbf{Noise and inexactness model} To emulate inexact gradient and operator evaluations that commonly arise in large-scale implementations, an epoch-dependent Gaussian perturbation $r_k\in\mathbb{R}^d$ with small standard deviation is added to all measurements at epoch $k$ and scaled by a multiplier $C\ge 0$, yielding epoch observations $\tilde{y}_i^{(k)}=y_i + C r_k$. This perturbation models both measurement noise and computation-induced errors and is reused consistently across methods where paired comparisons are required.

\textbf{Algorithms and hyperparameters} The proximal operator is implemented as soft-thresholding with regularization parameter $\lambda=10^{-5}$. Step-size is generally set to $\gamma=6.5\times10^{-8}$ to ensure stable proximal-gradient behavior, with the exception of the trajectory variance analysis (Figure~\ref{fig:trajectory_variance}) which utilizes $\gamma=10^{-3}$ to render update dynamics and angular deviations observable. EC-PRR integrates a velocity term with momentum coefficients studied in the set $\{0.0,0.5,0.7,0.9,0.99\}$ and the baseline PG-RR corresponds to momentum $=0$. All experiments use identical random seeds for data generation and permutation sequences to permit paired method comparisons. Epoch counts range up to 200 in convergence studies and 120 in robustness sweeps; metrics are recorded per epoch as specified in each experiment.

\textbf{Evaluation protocol and metrics} The primary performance metric is the normalized squared solution error $E_k=\|x_k - x^\star\|^2 / D$ recorded per epoch. To probe operator inexactness and update fidelity we compute epoch-wise averaged proximal-gradient residual norms, defined as the Euclidean norm between a per-sample update and the exact proximal-gradient mapping evaluated at the epoch start. We also analyze per-iteration stochastic gradient norm variance across epochs and the variance of inter-step angles between successive cumulative gradient sums inside each epoch; these latter metrics capture reshuffling-induced variability and stability of update directions.

\subsection*{Main Results}

\textbf{Momentum effect on convergence speed (Figure~\ref{fig:convergence_momentum})} Figure~\ref{fig:convergence_momentum} presents semilogarithmic convergence plots of the normalized squared solution error over 200 epochs for EC-PRR with different momentum coefficients and for the PG-RR baseline. The x-axis denotes epoch index and the y-axis shows the normalized squared error on a logarithmic scale. The curves compare momentum settings 0.0, 0.5, 0.7, 0.9, and 0.99 alongside the baseline. The plot demonstrates accelerated error decay for high momentum values, with momentum 0.9 and 0.99 producing the fastest convergence and the lowest steady-state errors. Lower or zero momentum yields a slow, nearly flat decay similar to PG-RR, underscoring that momentum is the critical component enabling EC-PRR to suppress error accumulation from reshuffling noise and inexact operators.

\begin{figure}[t]
  \centering
  \includegraphics[width=0.75\linewidth]{perspective/analysis_perspective_000/charts/convergence_momentum_comparison.png}
  \caption{Convergence curves illustrating the effect of different momentum parameters on the PG-RR algorithm. The horizontal axis represents the epoch number, ranging from 0 to 200, where each epoch corresponds to one pass over the measurement blocks. The vertical axis depicts the normalized squared error on a logarithmic scale, measuring the squared difference between the estimated and true vectors normalized by the number of samples. Curves correspond to the baseline PG-RR method without momentum (labeled as 'Baseline (PGRR)') and to variations of the algorithm with momentum parameters $\mu$ set to 0.00, 0.50, 0.70, 0.90, and 0.99, as indicated in the legend. Distinct marker styles and line patterns are used to differentiate between the momentum settings.}
  \label{fig:convergence_momentum}
\end{figure}

\textbf{Epoch-wise proximal-gradient residual norms (Figure~\ref{fig:epoch_residuals})} Figure~\ref{fig:epoch_residuals} shows epoch-wise averaged proximal-gradient residual norms for EC-PRR and PG-RR across 50 epochs. The residual norm quantifies the deviation between the per-sample update and the exact proximal-gradient mapping computed from the full gradient at the epoch start; the y-axis is logarithmic to reveal trends across orders of magnitude. The plot indicates that PG-RR attains numerically smaller instantaneous residuals than EC-PRR throughout epochs. However, interpreting these residuals requires algorithmic context: EC-PRR's residuals are larger because momentum modifies instantaneous proximal mappings by aggregating past gradients into a velocity term. The larger residuals therefore reflect a purposeful deviation from the instantaneous exact mapping in favour of error compensation and smoothing, which yields improved long-term convergence despite higher instantaneous discrepancy.

\begin{figure}[t]
  \centering
  \includegraphics[width=0.75\linewidth]{perspective/analysis_perspective_002/charts/epoch_residuals_comparison.png}
  \caption{Comparison of epoch-wise averaged proximal gradient residual norms between EC-PRR and PG-RR across 50 epochs. The horizontal axis shows epoch numbers and the vertical axis is a logarithmic scale of Euclidean norm residuals averaged over samples. Blue bars correspond to EC-PRR and orange bars to PG-RR. Larger residuals for EC-PRR reflect momentum-induced deviation from instantaneous proximal mappings.}
  \label{fig:epoch_residuals}
\end{figure}

\textbf{Robustness under varying observation noise (Figure~\ref{fig:robustness_noise})} Figure~\ref{fig:robustness_noise} compares convergence trajectories for EC-PRR and PG-RR at noise intensity levels $\{0,1,5,10,20\}$ over 120 epochs. The semilog plot overlays ten curves (two methods $\times$ five noise levels). EC-PRR consistently achieves lower normalized solution error across all noise intensities and exhibits smoother, steadier decay. As noise increases, both methods approach higher error floors, yet the gap between EC-PRR and PG-RR widens with noise intensity, indicating that momentum-based error compensation becomes increasingly advantageous in more corrupted scenarios.

\begin{figure}[t]
  \centering
  \includegraphics[width=0.75\linewidth]{perspective/analysis_perspective_004/charts/convergence_ecprr_vs_pgrr_noise.png}
  \caption{Convergence curves comparing EC-PRR and PG-RR across different noise levels. Horizontal axis is epoch number up to 120; vertical axis is normalized squared error on a logarithmic scale. Solid lines with circles denote EC-PRR; dashed lines with squares denote PG-RR. Colors identify noise magnitudes as specified in legend.}
  \label{fig:robustness_noise}
\end{figure}

\textbf{Variance of per-iteration stochastic gradient norms (Figure~\ref{fig:variance_grad_norms})} Figure~\ref{fig:variance_grad_norms} presents a boxplot comparison of the distribution of variances of per-iteration stochastic gradient norms across epochs for EC-PRR and PG-RR. Each variance value corresponds to an inner-iteration position (1..m) and is computed across epochs. The figure shows that EC-PRR displays larger raw variance in gradient norms relative to PG-RR. This observation should be interpreted in terms of accumulated gradient dynamics: momentum causes EC-PRR to aggregate gradients, which increases variability in instantaneous raw norms (computed prior to momentum aggregation) but produces smoother effective update directions when the velocity is taken into account. Hence, elevated raw variance does not contradict improved stability in iterate trajectories; rather, it reflects richer dynamics induced by momentum that are beneficial after aggregation.

\begin{figure}[t]
  \centering
  \includegraphics[width=0.75\linewidth]{perspective/analysis_perspective_005/charts/variance_comparison_gradient_norms.png}
  \caption{Boxplot comparing variance of per-iteration stochastic gradient norms over epochs for EC-PRR (momentum) and PG-RR (baseline). Horizontal axis denotes the method; vertical axis shows variance values of gradient norms computed at each inner iteration. The boxplots display median, IQR, whiskers, outliers, and mean (white circle). Color coding distinguishes methods.}
  \label{fig:variance_grad_norms}
\end{figure}

\textbf{Variance of accumulated gradient trajectory angles per epoch (Figure~\ref{fig:trajectory_variance})} Figure~\ref{fig:trajectory_variance} shows the mean and standard deviation across runs of the variance of inter-step angles between successive cumulative gradient sums for each epoch, using a step-size of $\gamma=10^{-3}$ to ensure observable geometric deviations. The x-axis is epoch index (1..50) and the y-axis is the variance of inter-step angles measured in radians squared. EC-PRR attains lower mean trajectory-angle variance and reduced run-to-run spread compared to PG-RR, indicating more coherent cumulative gradient directions within epochs. At the final epoch, the annotated values indicate approximately $1.13\times10^{-2}$ for EC-PRR versus $2.01\times10^{-2}$ for PG-RR, demonstrating that momentum reduces angular dispersion of the accumulated gradient trajectory and yields more stable update geometry.

\begin{figure}[t]
  \centering
  \includegraphics[width=0.75\linewidth]{perspective/analysis_perspective_011/charts/trajectory_variance_per_epoch.png}
  \caption{Line chart of per-epoch variance of accumulated gradient trajectory angles for EC-PRR (blue line with circles) and PG-RR (orange line with squares). Horizontal axis: epoch index 1--50; vertical axis: variance in radians squared. Shaded bands represent one standard deviation across runs. Annotated values near final epoch show mean variances for both methods.}
  \label{fig:trajectory_variance}
\end{figure}

\subsection*{Result Analysis}

\textbf{Validation of empirical claims against observed metrics} The convergence experiments corroborate the hypothesis that momentum-based error compensation accelerates convergence and reduces steady-state error in the presence of reshuffling and inexact operator computations. Figures~\ref{fig:convergence_momentum} and~\ref{fig:robustness_noise} together demonstrate that properly tuned momentum (notably in the vicinity of 0.9) materially improves epoch-wise error decay and yields superior robustness to observation noise. The residual-norm analysis in Figure~\ref{fig:epoch_residuals} validates that EC-PRR trades instantaneous proximal fidelity for aggregate stability: larger instantaneous residuals arise from the velocity-influenced updates, yet this mechanism prevents error amplification and supports improved asymptotic proximity to $x^\star$.

\textbf{Ablation-style comparison and parameter effect interpretation} By comparing EC-PRR runs across different momentum settings with the PG-RR baseline, the experiments function as an ablation on the momentum component. The momentum coefficient acts as the principal control knob for balancing error compensation and stability: zero momentum reduces EC-PRR to standard proximal RR behaviour; moderate values yield modest gains; high values ($\approx0.9$) offer the best empirical performance under the considered noise and step-size settings. The per-iteration variance results (Figure~\ref{fig:variance_grad_norms}) caution that raw stochastic gradient norms may appear more variable under EC-PRR; however, the trajectory-angle variance (Figure~\ref{fig:trajectory_variance}) confirms that velocity aggregation yields smoother effective update paths, thus reconciling the apparent paradox between raw-norm variance and update stability.

\textbf{Mechanistic interpretation based on method principles} The EC-PRR mechanism accumulates gradients into a velocity term $v_k=\beta v_{k-1}+\gamma g_k$, and applies a proximal mapping using $v_k$ rather than the instantaneous gradient. The experiments show that this accumulation diminishes the effective variance of update directions across inner iterations by averaging reshuffling noise, thereby reducing angular dispersion of cumulative gradients and improving the quality of descent directions. The residual-norm metric captures the deviation from the exact proximal-gradient mapping without velocity; EC-PRR intentionally incurs higher instantaneous residuals to realize lower long-term error through variance reduction and error compensation.

\textbf{Trade-offs between computational overhead and solution accuracy} The momentum update introduces negligible per-iteration overhead relative to the dominant matrix-vector operations in the measurement model, yet the empirical results indicate significant reductions in required epochs to reach a given error level. Thus, the additional computational cost of EC-PRR is justified by reduced overall runtime to a target accuracy in large-scale problems where per-epoch costs dominate.

\textbf{Implications for practical deployment and parameter tuning} The sensitivity studies and robustness sweeps indicate that practitioners should prioritize tuning the momentum coefficient when deploying EC-PRR in noisy or inexact environments. A momentum value near 0.9 proved effective across multiple settings; however, the optimal choice depends on the step-size, noise intensity $C$, and data conditioning. Paired experiments with identical reshuffling permutations and noise realizations are essential for reliable performance assessment and to prevent confounding factors in empirical comparisons.

\end{document}